\section{Background, Motivation and Problem}
\label{sec:motivation}
\subsection{What the decoder caches}
Let $X_\ell\in\mathbb{R}^{n\times d}$ be the normalized input of attention
layer $\ell$. In the full-MHA configuration studied here,
\begin{align}
 K_\ell &= \Rope(X_\ell W^K_\ell;P),\qquad
 V_\ell = X_\ell W^V_\ell, \label{eq:projections}\\
 A_{\ell,t} &= \softmax\!\left(
 Q_{\ell,t}K_{\ell,\leq t}^{\mathsf T}/\sqrt{d_h}+M_t\right),\nonumber\\
 O_{\ell,t} &= A_{\ell,t} V_{\ell,\leq t}. \label{eq:attention}
\end{align}
$P$ contains original logical positions, $M_t$ is the causal/validity mask,
$d_h$ is head dimension, and $W^K_\ell,W^V_\ell$ are frozen projections.
Rotary position embeddings (RoPE) make preserving positions part of the
representation contract, rather than a bookkeeping detail
\citep{su2021rope}.

For $L$ layers, batch size $B$, $H_{\mathrm{KV}}$ KV heads, sequence length
$n$ and $b$ bytes per scalar, a dense K/V pair occupies
\begin{equation}
 B_{\mathrm{dense}}(n)=2LBH_{\mathrm{KV}}nd_hb.
 \label{eq:densebytes}
\end{equation}
This expression is exact only for genuinely dense independent storage.
Views, shared buffers, padding and preallocation require backing-allocation
accounting. Our configuration has $L=32$, $H_{\mathrm{KV}}=32$, $d_h=128$,
$B=1$ and $b=2$. One processed token adds 0.50~MiB of dense KV. Producing
128 predictions from a prefilled prompt requires 127 further forwards;
thus FullKV grows by 63.50~MiB, not 64~MiB.

\subsection{Why source storage is attractive, and where it is not}
In full MHA, $H_{\mathrm{KV}}d_h=d$, so two projected arrays contain
$2nd$ scalars while their common input contains $nd$. The algebraic
opportunity precedes quantization: one source can supply both projections.
Quantizing that source offers additional storage reduction, at the price of
projection recomputation, quantization error and correction metadata.
There is no free latency benefit. A codec can save persistent bytes while
repeated unpacking, projection and synchronization cost more than ordinary
cache reads.

The dimensional argument changes under grouped-query attention (GQA):
$H_{\mathrm{KV}}d_h$ may be much smaller than $d$
\citep{ainslie2023gqa}. Storing the full normalized source can then be less
attractive. The current implementation explicitly restricts its geometry;
we do not extrapolate this study to Qwen, GQA, video or multiple images.

\subsection{Evidence loss and numerical error are different}
Token eviction removes a position from attention normalization. Source
quantization can instead retain a position with approximate K/V. Neither is
guaranteed to preserve the answer. Our protection policy assigns exact
storage to selected head-position regions while retaining the remaining
prompt positions through the source representation. A protected entry is
bitwise preserved; this does not imply that all evidence or downstream
logits are unchanged.

There is a second, distinct numerical issue: a custom reconstruction kernel
may differ from native projection of the \emph{same decoded source}. Our
construction checks and sparse numerical corrections address this transport
error. They do not remove the source quantization error relative to
unrestricted FullKV. Confusing those two reference boundaries would turn a
kernel-equivalence check into an unsupported answer-quality guarantee.

\subsection{The measured optimization problem}
We study a frozen request representation $R$ through four observables:
\begin{equation}
 \big(B_0(R),B_f(R),B_{\mathrm{dec,peak}}(R),T_{\mathrm{req}}(R)\big),
 \label{eq:observables}
\end{equation}
alongside corpus caption quality. $B_0$ and $B_f$ are initial and final
persistent decode-owned backing; $B_{\mathrm{dec,peak}}$ additionally
charges observed reconstruction/append overlap. It excludes vision/prefill
transients and is not a whole-request KV peak or process-HBM measurement.
The present study reports frozen operating points. It does not solve the
separate experiment of finding every method's smallest cache meeting a
common independently specified quality target.
